\documentclass[11pt]{amsart}

\usepackage[T1]{fontenc}
\usepackage{lmodern}
\usepackage{microtype}
\usepackage{mathtools,amssymb,amsthm}
\usepackage[hidelinks]{hyperref}

\hypersetup{
  pdftitle={A Steiner triple system of order 13 whose burning number increases
            on doubling and on tripling}
}

\newtheorem{theorem}{Theorem}
\newtheorem{proposition}[theorem]{Proposition}
\newtheorem{lemma}[theorem]{Lemma}

\newcommand{\STS}{\mathrm{STS}}
\newcommand{\PG}{\mathrm{PG}}
\newcommand{\AG}{\mathrm{AG}}
\newcommand{\bL}{b_{L}}

\title[An $\STS(13)$ whose burning number increases]
{A Steiner Triple System of Order~$13$ Whose Burning Number
 Increases on Doubling and on Tripling}

\date{}

\subjclass[2020]{05B07, 05C65, 05C57}
\keywords{Steiner triple system, burning number, doubling construction,
tripling construction, non-degenerate triple system, counterexample}

\begin{document}

\begin{abstract}
Burgess, Danziger, Jones, Marbach and Pike ask, in the Discussion of
\cite{BDJMP}, whether the projective and affine triple systems are the only
non-degenerate Steiner triple systems whose burning number increases on
doubling or tripling. We exhibit a counterexample. Let $H$ be the cyclic
$\STS(13)$ generated by the base blocks $\{0,1,4\}$ and $\{0,2,7\}$ modulo
$13$. Then $H$ is non-degenerate and, its order being neither of the form
$2^{n+1}-1$ nor of the form $3^{m}$, it is neither projective nor affine; yet
$b(H)=5$ while $b(D)=6$ for its double-and-add-one image $D$, an $\STS(27)$,
and $b(T)\ge6$ for its tripling image $T$, an $\STS(39)$. The values at orders
$13$ and $27$ are forced for \emph{every} Steiner triple system of those orders
by two unconditional bounds of \cite{BDJMP} itself, so no computation is
needed; they are also obtained directly from the definition of the burning
process. The refutation is thus arithmetic: it turns on the orders alone, and
it says nothing about orders at which those two bounds do not coincide.
\end{abstract}

\maketitle

\section{The question}

A \emph{Steiner triple system} of order $v$, written $\STS(v)$, is a pair
$(V,\mathcal B)$ with $|V|=v$ and $\mathcal B$ a set of $3$-subsets of $V$ such
that every unordered pair of points lies in exactly one block; such a system
exists precisely when $v\equiv1$ or $3\pmod 6$, and then
$|\mathcal B|=v(v-1)/6$.

\emph{Burning} a Steiner triple system \cite{BJP,BDJMP} is the following
one-player round-based process. Let $F(0)=\varnothing$. In round $r\ge1$ the
arsonist ignites one point $u_r\notin F(r-1)$, and \emph{simultaneously} every
point outside $F(r-1)$ that lies in a block whose other two points are both in
$F(r-1)$ catches fire; writing $\partial S$ for the set of such points,
\begin{equation}\label{eq:process}
  F(r)=F(r-1)\cup\partial F(r-1)\cup\{u_r\}.
\end{equation}
The propagation in \eqref{eq:process} is a \emph{single} step, not a closure.
The \emph{burning number} $b(H)$ is the least $k$ for which some choice of
sources $u_1,\dots,u_k$ gives $F(k)=V$. In the \emph{lazy} variant the fire
spreads to a fixpoint and no source is added after the first round, so the lazy
burning number $\bL(H)$ is the least size of a set of points whose closure
under ``the third point of the block through a pair'' is all of $V$. A
subsystem of $H$ is a subset of $V$ closed under that operation. Following
\cite{BDJMP}, $H$ is \emph{non-degenerate} if every maximal proper
subsystem $S$ of $H$ satisfies $\bL(S)=\bL(H)-1$.

Two classical recursions appear in the question. \emph{Doubling}
(double-and-add-one, \cite[Construction~2.15]{Handbook}) takes an $\STS(v)$ on
$V$, a set $W$ of $v+1$ further points and a $1$-factorization
$F_0,\dots,F_{v-1}$ of the complete graph on $W$, and produces the $\STS(2v+1)$
on $V\cup W$ whose blocks are those of the original together with
$\{x\}\cup e$ for every $x\in V$ and every $e\in F_x$. \emph{Tripling}
(\cite[Construction~2.17]{Handbook}) produces an $\STS(3v)$ on $V\times\mathbb
Z_3$ whose blocks are the $v$ vertical triples
$\{(x,0),(x,1),(x,2)\}$ together with
$\{(x,i),(y,j),(z,k)\}$ for every block $\{x,y,z\}$ of the original and every
$(i,j,k)\in\mathbb Z_3^3$ with $i+j+k\equiv0$. The projective triple systems
$\PG(n,2)$, of order $2^{n+1}-1$, are exactly those obtained by repeated
doubling from $\STS(3)$; the affine triple systems $\AG(m,3)$, of order $3^m$,
are exactly those obtained by repeated tripling from $\STS(3)$.

The question refuted here stands in the Discussion of \cite{BDJMP} and carries
no number there, so we quote it from the e-print source together with the
sentence before it, the two internal cross-references redacted:

\begin{footnotesize}
\begin{verbatim}
Theorems [...] and [...] also show that doubling a
projective triple system or tripling an affine triple system
increases the burning number by one.  We ask if these are the
only non-degenerate systems whose burning numbers increase on
doubling or tripling.
\end{verbatim}
\end{footnotesize}

\noindent
We read ``increase on doubling \emph{or} tripling'' in its union sense, which
is the weaker target.

\begin{theorem}\label{thm:main}
Let $H$ be the $\STS(13)$ on $\mathbb Z_{13}$ whose $26$ blocks are listed in
\eqref{eq:H}, let $D$ be the $\STS(27)$ with block set \eqref{eq:D}, obtained
from $H$ by doubling, and let $T$ be the $\STS(39)$ obtained from $H$ by
tripling. Then
\begin{enumerate}
\item[(a)] $H$ is non-degenerate;
\item[(b)] $13\ne2^{n+1}-1$ for every $n\ge0$ and $13\ne3^{m}$ for every
      $m\ge0$, so $H$ is neither a projective nor an affine triple system;
\item[(c)] $b(H)=5$, $b(D)=6$ and $b(T)\ge6$.
\end{enumerate}
So $H$ is a non-degenerate Steiner triple system, neither projective nor
affine, whose burning number increases on doubling and increases on tripling.
With ``or'' read as a union, the quoted question therefore has answer no.
\end{theorem}

\section{The witness}

Let $H$ be the cyclic $\STS(13)$ on $\mathbb Z_{13}$ obtained by developing the
base blocks $\{0,1,4\}$ and $\{0,2,7\}$ modulo $13$. Its $26$ blocks are
\begin{equation}\label{eq:H}
\begin{aligned}
&\{0,1,4\},\{0,2,7\},\{0,3,12\},\{0,5,11\},\{0,6,8\},\{0,9,10\},\\
&\{1,2,5\},\{1,3,8\},\{1,6,12\},\{1,7,9\},\{1,10,11\},\\
&\{2,3,6\},\{2,4,9\},\{2,8,10\},\{2,11,12\},\\
&\{3,4,7\},\{3,5,10\},\{3,9,11\},\\
&\{4,5,8\},\{4,6,11\},\{4,10,12\},\\
&\{5,6,9\},\{5,7,12\},\\
&\{6,7,10\},\{7,8,11\},\{8,9,12\}.
\end{aligned}
\end{equation}
These are $26=13\cdot12/6$ triples covering each of the $78$ pairs of
$\mathbb Z_{13}$ exactly once.

For the doubling, let $W=\{w_0,\dots,w_{12},w_\infty\}$ and take the standard
$1$-factorization of the complete graph on $W$ indexed by $\mathbb Z_{13}$,
\[
  F_i=\bigl\{\{w_\infty,w_i\}\bigr\}\cup
      \bigl\{\{w_{i+j},w_{i-j}\}:j=1,\dots,6\bigr\}
  \qquad(i\in\mathbb Z_{13}),
\]
subscripts modulo $13$. Each $F_i$ is a perfect matching of the $14$ points of
$W$, and the thirteen $F_i$ partition all $91$ edges, because
$\{w_a,w_b\}$ with $a\ne b$ lies in $F_i$ for the unique $i=(a+b)/2$ ($2$ being
invertible modulo $13$) and $\{w_\infty,w_a\}$ lies in $F_a$ alone. Let $D$ be
the $\STS(27)$ on $\mathbb Z_{13}\cup W$ whose block set is
\begin{equation}\label{eq:D}
  \eqref{eq:H}\ \cup\
  \bigl\{\{i\}\cup e\ :\ i\in\mathbb Z_{13},\ e\in F_i\bigr\},
\end{equation}
that is, the $26$ blocks of \eqref{eq:H} together with the $91$ blocks
$\{i,w_\infty,w_i\}$ and $\{i,w_{i+j},w_{i-j}\}$ for $i\in\mathbb Z_{13}$ and
$j=1,\dots,6$; in all $26+91=117=27\cdot26/6$ blocks. Finally $T$ is the
$\STS(39)$ on $\mathbb Z_{13}\times\mathbb Z_3$ given by the tripling rule of
\S1 applied to \eqref{eq:H}, with $13+9\cdot26=247=39\cdot38/6$ blocks.

\section{Proof of Theorem~\ref{thm:main}}

\subsection*{(a) $H$ is non-degenerate}
A proper subsystem $S$ of an $\STS(v)$ has order $w$ with $v\ge2w+1$: fix a
point $x\notin S$ and send $p\in S$ to the third point $f(p)$ of the block
$\{x,p,f(p)\}$. If $f(p)=f(q)$ for $p\ne q$ then the block through $f(p)$ and
$x$ would meet $S$ twice, forcing $x\in S$; and $f(p)\in S$ would likewise force
$x\in S$. So $f$ is an injection of $S$ into $V\setminus(S\cup\{x\})$, whence
$w\le v-w-1$. For $v=13$ this gives $w\le6$, so the only admissible orders of a
proper subsystem are $w=1$ and $w=3$: $H$ contains no subsystem of order $7$ or
$9$. The subsystems of order $3$ are exactly the blocks, and each block is a
maximal proper subsystem. A block, as an $\STS(3)$, is generated by any two of
its points and by no single point, so $\bL=2$ there. On the other hand no pair
of points generates $H$, since the closure of a pair is the block through it,
while $H$ \emph{is} generated by three points: the closure of $\{0,1,2\}$ already
contains $\{0,1,2,4,5,7\}$, the blocks through the three pairs being
$\{0,1,4\}$, $\{1,2,5\}$ and $\{0,2,7\}$, and a subsystem of order at least $6$
must have order $13$, the admissible orders being only $1$, $3$ and $13$. Hence
$\bL(H)=3$ and every maximal proper subsystem $S$ has
$\bL(S)=2=\bL(H)-1$: $H$ is non-degenerate.

\subsection*{(b) $H$ is neither projective nor affine}
$2^{n+1}-1$ takes the values $7$ and $15$ on either side of $13$, and $3^m$
takes the values $9$ and $27$. Both are statements about the \emph{order}, so
they hold for every $\STS(13)$ and require no isomorphism test.

\subsection*{(c) The burning numbers}
Define, with $\binom{\cdot}{2}$ the binomial coefficient,
\begin{equation}\label{eq:h}
  h(1)=1,\qquad h(r)=\binom{h(r-1)}{2}-2h(r-1)+3r-2 ,
\end{equation}
so that $h(1),\dots,h(6)=1,\ 2,\ 4,\ 8,\ 25,\ 266$. Write $\rho_v$ for the
least $r$ with $v\le h(r)$, and $U(v)=\lceil\log_2v\rceil+1$. Two theorems of
\cite{BDJMP} are unconditional and hold for every Steiner triple system:
\begin{equation}\label{eq:sandwich}
  \rho_v\le b(K)\le U(v)\qquad\text{for every }\STS(v)\ K .
\end{equation}
The lower bound is the theorem of \cite{BDJMP} whose engine is the statement
that at most $h(r)$ points can be alight at the end of round $r$; the upper
bound is the bound $b\le\lceil\log_2v\rceil+1$ proved there for all Steiner
triple systems. Now
\[
  \begin{array}{lll}
    h(4)=8<13\le25=h(5) &\Rightarrow \rho_{13}=5, & U(13)=4+1=5;\\
    h(5)=25<27\le266=h(6) &\Rightarrow \rho_{27}=6, & U(27)=5+1=6;\\
    h(5)=25<39\le266=h(6) &\Rightarrow \rho_{39}=6, & U(39)=6+1=7.
  \end{array}
\]
In the first two lines the two bounds of \eqref{eq:sandwich} coincide, so
$b(K)=5$ for \emph{every} $\STS(13)$ and $b(K)=6$ for \emph{every} $\STS(27)$.
In particular $b(H)=5$ and $b(D)=6$ whatever $1$-factorization is used in
\eqref{eq:D}, so the ``for some image'' and ``for every image'' readings of
``increases on doubling'' agree here. Finally $b(T)\ge\rho_{39}=6>5=b(H)$.
This settles (c), and with it Theorem~\ref{thm:main}, by counting alone.

\section{What is not settled}

The two equalities used in (c) hold for every system of order $13$ and every
system of order $27$; they are facts about those orders and not about $H$. So
the witness contradicts the question as printed without bearing on the
structural phenomenon that \cite{BDJMP} exhibits at the projective and affine
systems, and nothing here shows that a projective/affine characterisation
fails at orders where the two bounds of \eqref{eq:sandwich} do not coincide.
A referee may reasonably read Theorem~\ref{thm:main} as a small-order artefact
of those bounds.

The increase from $H$ to $D$ is by exactly one, both endpoints being pinned.
Whether the increase from $H$ to $T$ is by exactly one is not determined here:
\eqref{eq:sandwich} confines $b(T)$ to $\{6,7\}$. The reading on which the
question asks for an increase \emph{by one} on doubling \emph{and} on tripling
is therefore not settled by this witness.

\section{Verification}

The program accompanying this note reads the $26$ blocks of \eqref{eq:H} as
they are printed above and re-derives the quantities asserted here: that
\eqref{eq:H} is an $\STS(13)$ and coincides with the cyclic development of
$\{0,1,4\}$ and $\{0,2,7\}$; that the $F_i$ are a $1$-factorization and that
$D$ and $T$ are Steiner triple systems of orders $27$ and $39$ restricting to
$H$ as claimed; that $H$ has exactly $41$ subsystems, that its $26$ maximal
proper subsystems are its blocks, and that $\bL(H)=3$ with $\bL=2$ on each of
them, so that $H$ is non-degenerate; the recursion \eqref{eq:h} and the values
$\rho_{13},\rho_{27},\rho_{39}$ and $U(13),U(27),U(39)$, with
$\lceil\log_2v\rceil$ evaluated exactly by bit length rather than by a
floating-point logarithm. It also recomputes $b(H)=5$, $b(D)=6$ and $b(T)>5$
from the definition \eqref{eq:process} alone, by enumerating every burnt set
reachable at the end of each round, so that (c) does not rest on
\eqref{eq:sandwich}; the sources $0,1,2,3,6$ burn $H$ in five rounds. As
controls it reproduces $b(\STS(3))=3$, $b(\PG(2,2))=4$, $b(\AG(2,3))=5$,
$b(\PG(3,2))=5$ and $b(\AG(3,3))=6$, all values published in \cite{BDJMP}, and
checks that the tripling rule carries $\AG(2,3)$ to $\AG(3,3)$. It is
Python~3.9 or later, standard library only, exact integer arithmetic
throughout, and it runs in a few seconds. Its transcript accompanies it.

Its scope has limits, which it also prints. It does not fetch or parse the
e-print of \cite{BDJMP}: the quoted question is transcribed from the source,
and only the values of $h$ are recomputed. It does not consult \cite{Handbook};
the doubling and tripling rules are implemented from the statements reproduced
in \S1. No isomorphism census is performed. The exact value of $b(T)$, which is
$6$ or $7$, is not determined.

\begin{thebibliography}{9}

\bibitem{BDJMP}
A.~C. Burgess, P.~H. Danziger, C.~W. Jones, T.~G. Marbach, and D.~A. Pike,
\emph{Burning Steiner triple systems},
arXiv:2608.25627v1, 2026.
\href{https://arxiv.org/abs/2608.25627}{arXiv:2608.25627}.
The question refuted here is unnumbered in that e-print; it is quoted in full
in \S1.

\bibitem{BJP}
A.~C. Burgess, C.~W. Jones, and D.~A. Pike,
\emph{Extending graph burning to hypergraphs},
Ars Math. Contemp. \textbf{26} (2026), Paper P3.07.
\href{https://arxiv.org/abs/2403.01001}{arXiv:2403.01001}.

\bibitem{Handbook}
C.~J. Colbourn and J.~H. Dinitz (eds.),
\emph{Handbook of Combinatorial Designs}, 2nd ed.,
Chapman \& Hall/CRC, Boca Raton, 2007.
Constructions~2.15 (double-and-add-one) and~2.17 (tripling). Not consulted
directly; the two constructions are used here in the standard form reproduced
in \S1.

\end{thebibliography}

\end{document}
